% ===================== ALCANCE =====================
\section{Alcance y Limitaciones}
\label{sec:alcance}

\subsection{Alcance del Estudio}

\begin{table}[htbp]
    \centering
    \caption{Elementos incluidos en el alcance del an\'{a}lisis.}
    \label{tab:alcance_incluido}
    \setcounter{itemcount}{0}
    \begin{tabularx}{\textwidth}{@{}NX@{}}
        \toprule
        \multicolumn{1}{c}{\textbf{Item}} & \textbf{Descripcion} \\
        \midrule
        & Modelo de elementos finitos de la l\'{i}nea SIM-015 (job CAESAR P2603-PR-SIM-015), tuber\'{i}a del filtro de fibras PP1, generado desde el isom\'{e}trico PCF114. \\
        & Evaluaci\'{o}n de esfuerzos seg\'{u}n ASME B31.3-2016 para nueve casos de carga: dos de operaci\'{o}n (OPE), dos sostenidos alternos (Alt-SUS), dos sostenidos (SUS) y tres de expansi\'{o}n (EXP). \\
        & Reporte de desplazamientos nodales, esfuerzos por elemento y cargas en soportes para los casos de carga relevantes. \\
        & Verificaci\'{o}n de cumplimiento de c\'{o}digo (\textit{code compliance}) y elaboraci\'{o}n del presente informe t\'{e}cnico con c\'{o}digo de documento I24.104-PP30-PP01-P-DOCS-190. \\
        \bottomrule
    \end{tabularx}
\end{table}

\subsection{Limitaciones del Estudio}

\begin{table}[htbp]
    \centering
    \caption{Elementos excluidos del alcance del an\'{a}lisis.}
    \label{tab:alcance_excluido}
    \setcounter{itemcount}{0}
    \begin{tabularx}{\textwidth}{@{}NX@{}}
        \toprule
        \multicolumn{1}{c}{\textbf{Item}} & \textbf{Descripcion} \\
        \midrule
        & Cargas ocasionales de viento y sismo: no reportadas en el modelo; su evaluaci\'{o}n, de ser requerida, corresponde a un an\'{a}lisis posterior. \\
        & Cargas en boquillas de equipos frente a criterios de fabricante: la comparaci\'{o}n de las cargas en las conexiones con los equipos (anclas de los nodos 140, 320 y 510) con los admisibles de cada equipo requiere los datos del fabricante, no disponibles a la fecha. \\
        & Fatiga detallada por ciclos de arranque/parada: el criterio de expansi\'{o}n de B31.3 con factor $f$ cubre el rango t\'{e}rmico; no se ejecut\'{o} an\'{a}lisis de fatiga espec\'{i}fico. \\
        & Verificaci\'{o}n en campo de la soporter\'{i}a: las propiedades de los soportes corresponden a las definidas en el isom\'{e}trico de dise\~{n}o. \\
        \bottomrule
    \end{tabularx}
\end{table}

\subsection{Advertencias del Modelo}

La corrida de CAESAR II report\'{o} advertencias de modelado que se declaran a
continuaci\'{o}n. Algunos componentes del archivo de intercambio no tra\'{i}an
material asignado en la especificaci\'{o}n y recibieron el material por defecto
durante la importaci\'{o}n; el material de evaluaci\'{o}n de todo el modelo se
fij\'{o} posteriormente en ASTM A312 TP 304L, consistente con el admisible
reportado de \num{115142.4}~\si{kPa} (secci\'{o}n~\ref{sec:resultados}).
Adem\'{a}s, los soportes sin tipo definido en el isom\'{e}trico se importaron
como soportes verticales, consistentes con la soporter\'{i}a de la
Figura~\ref{fig:nodos-soporte}. Estas advertencias no afectan el veredicto de
cumplimiento de c\'{o}digo, pero deben considerarse ante cualquier
modificaci\'{o}n futura del modelo.
